\documentclass[10pt,a4paper]{amsart}
\usepackage[margin=19mm]{geometry}
\usepackage{amsmath,amssymb,mathtools}
\usepackage[scr=rsfs,cal=euler]{mathalfa}
\usepackage{xfrac}
\usepackage[unicode,hidelinks,bookmarks=false]{hyperref}
\newcommand{\ie}{i.e.\ }
\newcommand{\cf}{cf.\ }
\newcommand{\resp}{respectively\ }
\newcommand{\Subsection}{\subsection}
\providecommand{\nobreakdash}{\nobreak\hbox{-}}
\setlength{\emergencystretch}{2em}
\multlinegap0pt
\usepackage{mathtools}
\usepackage[shortlabels]{enumitem}
\usepackage{xcolor}
\usepackage{amssymb}
\newtheorem{lem}{Lemma}
\newtheorem{remark}{Remark}
\usepackage{cleveref}
\crefname{lem}{Lemma}{Lemmas}
\crefname{remark}{Remark}{Remarks}
\newcommand{\N}{\mathbb{N}}
\newcommand{\R}{\mathbb{R}}
\newcommand{\de}{\mathrm{d}}
\newcommand{\eps}{\varepsilon}
\title{Birkhoff genericity on affine subspaces in horospheres}
\author{Nimish A. Shah, Pengyu Yang}
\date{Source: arXiv:2606.08199v1 (CC BY 4.0)}
\begin{document}
\maketitle

\noindent\emph{Source and licence.} Birkhoff genericity on affine subspaces in horospheres. Version arXiv:2606.08199v1 (CC BY 4.0), distributed by arXiv under the Creative Commons Attribution 4.0 International licence, http://creativecommons.org/licenses/by/4.0/, as recorded in the arXiv licence metadata for this exact version and shown on the abstract page https://arxiv.org/abs/2606.08199v1. Full licence notice: \texttt{licenses/arxiv-2606.08199-CC-BY-4.0.txt}.\par\smallskip

\noindent\textit{Editorial scope.} Counted unit: the article's lem compact (source0.tex lines 1314--1324). The article's complete original proof of it is delivered with it (source0.tex lines 1325--1372, one \texttt{proof} environment of 48 lines). No mathematics is written, repaired or re-derived: the statement and the proof are the article's own lines, in the article's own notation, and no retained line is substituted or re-flowed.  Retained uncounted as the source-local context the proof needs: lines 913--919.  Nothing was omitted from any retained span, so the package carries no omission record.  No retained line contains a citation, so the wrapper carries no bibliography and no bibliography entry.  No retained line contains a figure environment, an image-inclusion command or drawing code, so no image is delivered and the package declares no asset dependency.  Editorial formatting only, in addition: an AMS \texttt{amsart} preamble that restates the article's own declarations and macros used by the retained lines, namely the environments \texttt{lem}, \texttt{remark} on separate counters, together with the standard packages those lines need.  The retained environments print in this document as Lemma 1, Remark 1 in order of appearance, whereas the article prints the same items with the numbering of its own full document.  The complete native source of the article as distributed in the arXiv e-print for this exact version is bundled unchanged as \texttt{sources/202314/source0.tex}.
	\begin{lem}\label{lem:contraction_alpha}
		Let $0<\theta<\frac{d}{d+1}$ and $\eps>0$. There exists $C\geq 1$ such that for any $t\geq0$ there exists $b\geq 0$ such that for any $x\in X$ we have
		\begin{equation}\label{eq:contraction hypothesis alphaEpsTheta}
		\int_{I^d}\alpha_\eps^{\theta}(c_tu(s)x)\de s\leq
		Ce^{-\theta t}\alpha_\eps^{\theta}(x)+b.
		\end{equation}
	\end{lem}

\begin{lem} \label{lem:back to compact}
        Let $\sigma_0>1$ be given. Then there exist $t_1\geq 1$ and $l_1\geq 1$ such that for any $l\geq l_1$, and any $x\in X'$ the following holds: There exist 
        \[
        \tau\leq t_1+\sigma_0\theta^{-1}\log^+ (2l^{-1}\alpha_\eps^\theta(x)) \text{ and } s\in [-1,1]^d
        \]
        such that
        \[
        c_{\tau}u(s)x \in X_l.  
        \]
        Here $\log^+(y):=\max\{ \log y,0 \}$ for every $y>0$.
    \end{lem}

    \begin{proof}
       Let $C>1$ be such that the conclusion of \Cref{lem:contraction_alpha} holds. Take $t_1\geq 1$ sufficiently large such that 
        \[
        \frac{\theta}{\theta-t_1^{-1}\log C}\leq\sigma_0.
        \]
        We will write $t=t_1$ for simplicity. By \Cref{lem:contraction_alpha}, there exists $b>0$ such that 
        \[
        \int_{I^d}\alpha_\eps^{\theta}(c_tu(s)x)\de s\leq
		Ce^{-\theta t}\alpha_\eps^{\theta}(x)+b.
        \]
        Iterating $n$ times,
        \begin{equation} \label{eq:exponential contraction}
        \begin{split}
            \int_{\R^d}\alpha_\eps^{\theta}(c_{nt}u(s)x)\de \nu^{(n)}(s)
            &\leq C^ne^{-n\theta t}\alpha_\eps^{\theta}(x)+\frac{1-C^ne^{-n\theta t}}{1-Ce^{-\theta t}}b \\
            &< C^ne^{-n\theta t}\alpha_\eps^{\theta}(x)+\frac{b}{1-Ce^{-\theta t}},
        \end{split}    
        \end{equation}
        where $\nu^{(n)}$ is a probability measure on $\R^d$ supported on $(1-e^{-t})^{-1}I^d$. Indeed, $\nu^{(1)}$ is the Lebesgue measure restricted to $I^d=[-1/2,1/2]^d$, and for each $n\geq 2$,
        \begin{align*}
&\int_{I^d}\left(\int_{\R^d}\alpha_\eps^\theta(c_{(n-1)t}u(s_1)c_{t}u(s)x)\,d\nu^{(n-1)}(s_1)\right)\,d\nu^{(1)}(s)\\
&=\int_{I^d}\left(\int_{\R^d}\alpha_\eps^\theta(c_{nt}u(e^{-t}s_1+s)x)\,d\nu^{(n-1)}(s_1)\right)\,d\nu^{(1)}(s)\\
&=\int_{\R^d}\alpha_\eps^{\theta}(c_{nt}u(s)x)\,d\nu^{(n)}(s),
\end{align*}
    where $\nu^{(n)}$ is the convolution of $(e^{-t}I_d)_\ast \nu^{(n-1)}$ and $\nu^{(1)}$, where $e_{-t}I_d$ denotes the multiplication by $e^{-t}$ on $\R^d$. 

    \begin{remark} \label{rem:nu*n}
       For later purpose we note that $\de\nu^{(n)}=f_n\de s$, where $f_n\geq 0$, $f_n$ is symmetric about each coordinate axis, $\int f_n=1$, and if $t_1\geq 2$, then $f_n\geq 4^{-d}\mathbf{1}_{I^d}$; that is, $\nu^{(1)}\leq 4^d\nu^{(n)}$ for all $n$.  
    \end{remark}
        
        Take $l_1=\frac{2b}{1-Ce^{-\theta t}}$. Let $l\geq l_1$ and choose smallest $n\in\N$ such that 
        \[
        C^ne^{-n\theta t}\alpha_\eps^{\theta}(x)\leq l/2. 
        \]
        So, we choose
        $n=\left\lceil \frac{\log^+ (2l^{-1}\alpha_\eps^\theta(x))}{\theta t-\log C} \right\rceil$.
        Then, the right hand side of \eqref{eq:exponential contraction} is at most $l/2+l_1/2\leq l$. Hence, by \eqref{eq:exponential contraction}, there exist 
        \[
        s\in {(1-e^{-t})}^{-1}I^d\subset [-1,1]^d
        \]
        such that $\alpha_\eps^{\theta}(c_{nt}u(s)x)\leq l$.  Now, $t=t_1$ and
        \begin{align*}
            nt&\leq t+\frac{t}{\theta t-\log C}\cdot \log^+ (2l^{-1}\alpha_\eps^\theta(x)) \\
            &=t_1+\theta^{-1}\cdot \frac{\theta}{\theta-t_1^{-1}\log C}\cdot \log^+ (2l^{-1}\alpha_\eps^\theta(x))\\
            &\leq t_1+\theta^{-1}\sigma_0\log^+ (2l^{-1}\alpha_\eps^\theta(x)).
        \end{align*}
       We conclude the proof of the lemma by choosing $\tau=nt_1$.
    \end{proof}

\end{document}
